\chapter{Hip\'oteses do Estudo}
\label{cap:hipoteses}

\textbf{H1} -- A arquitetura proposta gerar\'a orienta\c{c}\~oes em renda
fixa com precis\~ao num\'erica verific\'avel (c\'alculos de rentabilidade
l\'iquida, impacto de IR/IOF e duration) e conformidade regulat\'oria
comprovada, demonstrando viabilidade como ferramenta de assessoria
automatizada em ambiente regulado.

\textit{Fundamenta\c{c}\~ao}: estudos sobre rob\^os-consultores
\autocite{dacunto2019} documentam que modelos algor\'itmicos reduzem os
vieses comportamentais humanos de tomada de decis\~ao. A segrega\c{c}\~ao
entre c\'alculo composto determin\'istico (Python) e gera\c{c}\~ao de
linguagem (LLM) elimina a principal fonte de erro em modelos generativos
puros -- as alucina\c{c}\~oes matem\'aticas --, conferindo auditabilidade
integral \`as sa\'idas. A precis\~ao num\'erica \'e aferida frente ao
padr\~ao de refer\^encia (\textit{ground truth}) independentemente
verificado nos simuladores oficiais (Tesouro Transparente, calculadoras
ANBIMA), conforme detalhado no objetivo espec\'ifico c, dispensando a
necessidade de um comparador externo.

\textbf{H2} -- A avalia\c{c}\~ao das recomenda\c{c}\~oes do agente por
gerentes de front-office do BRB -- aplicada via question\'ario System
Usability Scale (SUS) nos crit\'erios de clareza, precis\~ao e
adequa\c{c}\~ao regulat\'oria -- atingir\'a score SUS igual ou superior a
68 pontos, classificado na literatura como patamar ``Aceit\'avel/Bom''
\autocite{bangor2008,sauro2011}.

\textit{Fundamenta\c{c}\~ao}: a hip\'otese apoia-se na identifica\c{c}\~ao
de lacunas na qualidade e profundidade da assessoria nos canais f\'isicos
de varejo, frequentemente decorrentes da sobrecarga operacional dos
gerentes comerciais. Um agente de software especializado, com acesso
instant\^aneo a todas as curvas de juros e resolu\c{c}\~oes aplic\'aveis,
tende a produzir recomenda\c{c}\~oes mais robustas para carteiras padr\~ao.
O score SUS de 68 pontos (escala de 0 a 100) \'e o limiar emp\'irico
estabelecido na literatura como patamar m\'inimo de aceitabilidade de
sistemas interativos, calculado pela metodologia padr\~ao do question\'ario
(soma das pontua\c{c}\~oes convertidas dos 10 itens). O instrumento ser\'a
aplicado em sua forma original, sem convers\~ao para escalas alternativas.

\textbf{H3} -- O framework desenvolvido no BRB apresentar\'a arquitetura
modular suficientemente desacoplada de especificidades institucionais
para ser adaptado a outros bancos p\'ublicos estaduais, com esfor\c{c}o de
customiza\c{c}\~ao restrito aos m\'odulos de integra\c{c}\~ao de dados
internos e parametriza\c{c}\~ao regulat\'oria local, indicando potencial
de escala para pol\'iticas p\'ublicas de inclus\~ao financeira no \^ambito
da ENEF.

\textit{Fundamenta\c{c}\~ao}: a modularidade da arquitetura, apoiada em
APIs p\'ublicas (BCB/ANBIMA) e no mapeamento de demandas de varejo via
bases p\'ublicas online, reduz significativamente as amarras sist\^emicas
corporativas da aplica\c{c}\~ao, facilitando a transposi\c{c}\~ao do
modelo para estruturas banc\'arias cong\^eneres e sua eventual
incorpora\c{c}\~ao como diretriz de pol\'itica p\'ublica.
